% part: intuitionistic-logic
% chapter: propositions-as-types
% section: terms-to-proofs

\documentclass[../../../include/open-logic-section]{subfiles}

\begin{document}

\olfileid{int}{pty}{tp}

\olsection{સાબિતી-પદોમાંથી \usetoken{P}{derivation} પુનઃપ્રાપ્ત કરવી}

હવે બીજી દિશા વિચારીએ: સાબિતી-પદ~$N$માંથી
\Log{N2i}ની !!a{proof} પાછી મેળવવી.
સામાન્ય રીતે આ ફક્ત એવા સંદર્ભની સાપેક્ષે શક્ય છે,
જેમાં આપેલા સાબિતી-પદમાં મુક્ત રીતે આવતાં દરેક ચલ~$x$
માટે $x : !A$ની જોડી હોય.
પરંતુ મુક્ત ચલ વિનાના ઉદાહરણથી શરૂ કરીએ, એટલે આ પદથી:
\[
\lambd[\typeof{x}{!A \land
  !B}][\pair{\proj{2}{x}}{\proj{1}{x}}]
\]
તેનું સ્વરૂપ $\lambd[\typeof{x}{!A \land !B}][N_1]$ છે.
\Log{tN2i}માં આ સ્વરૂપનું પદ આપતો એકમાત્ર નિયમ
\Intro{\lif} છે; એટલે $N$માં સંકેતિત !!{proof}નો
અંત~\Intro{\lif}થી જ થાય છે, એટલે તે આમ પૂરો થાય છે:
  \[
  \Axiom$x : !A\land !B \fCenter \pair{\proj{2}{x}}{\proj{1}{x}} : 
    !C$
  \RightLabel{\Intro\lif}
  \UnaryInf$\fCenter \lambd[\typeof{x}{!A \land
  !B}][\pair{\proj{2}{x}}{\proj{1}{x}}] : (!A \land !B) \lif !C$
  \DisplayProof
  \]
હજુ $!C$ શું છે તે જાણતા નથી, પરંતુ પૂર્વાંગ
$!A \land !B$ છે અને આધારના સંદર્ભમાં
$x:!A \land !B$ આવવું જ જોઈએ તે જાણીએ છીએ.

હવે આધાર સાથે સંકળાયેલું સાબિતી-પદ નક્કી થઈ ચૂક્યું છે:
$\pair{\proj{2}{x}}{\proj{1}{x}}$.
\emph{તેમાં} સંકેતિત !!{proof}નો અંત \Intro{\land}થી જ થાય.
હજુ $!C$ શું છે તે જાણતા નથી, પરંતુ
તે~\Intro{\land}નું નિષ્કર્ષ હોવાથી સંયોજન જ હોય,
એટલે $!C \ident (!C_1 \land !C_2)$.
પુનઃરચના આગળ વધારીએ:
  \[
  \Axiom$x : !A\land !B \fCenter \proj{2}{x} : !C_1$
  \Axiom$x : !A\land !B \fCenter \proj{1}{x} : !C_2$
  \RightLabel{\Intro\land}
  \BinaryInf$x : !A\land !B \fCenter \pair{\proj{2}{x}}{\proj{1}{x}} : 
    !C_1 \land !C_2$
  \RightLabel{\Intro\lif}
  \UnaryInf$\fCenter \lambd[\typeof{x}{!A \land
  !B}][\pair{\proj{2}{x}}{\proj{1}{x}}] : (!A \land !B) \lif (!C_1 \land !C_2)$
  \DisplayProof
  \]
સાબિતી-પદ $\pi_2(x)$ સૂચવે છે કે ડાબી તરફનો સૌથી ઉપરનો
સિક્વન્ટ \Elim{\land}[2]થી મળ્યો હતો, એટલે આધારનું
સ્વરૂપ $!D \land !C_1$ છે.
જમણી તરફ $!C_2$ આપતું અનુમાન \Elim{\land}[1] જ હોય
અને આધારનું સ્વરૂપ $!C_2 \land !E$ હોય તે નક્કી કરી શકીએ.
  \[
  \Axiom$x : !A\land !B \fCenter x: !D\land !C_1$
  \RightLabel{\Elim\land[2]}
  \UnaryInf$x : !A\land !B \fCenter \proj{2}{x} : !C_1$
  \Axiom$x : !A\land !B \fCenter x: !C_2\land !E$
  \RightLabel{\Elim\land[1]}
  \UnaryInf$x : !A\land !B \fCenter \proj{1}{x} : !C_2$
  \RightLabel{\Intro\land}
  \BinaryInf$x : !A\land !B \fCenter \pair{\proj{2}{x}}{\proj{1}{x}} : 
    !C_1 \land !C_2$
  \RightLabel{\Intro\lif}
  \UnaryInf$\fCenter \lambd[\typeof{x}{!A \land
  !B}][\pair{\proj{2}{x}}{\proj{1}{x}}] : (!A \land !B) \lif (!C_1 \land !C_2)$
  \DisplayProof
  \]
સૌથી ઉપરના સિક્વન્ટોનાં સાબિતી-પદો હવે ચલ છે,
એટલે તેઓ સ્વયંસિદ્ધો જ હોવા જોઈએ.
આથી $!C_1$, $!C_2$, $!D$ અને $!E$ નક્કી થાય છે:
$!D \ident !A$ અને $!C_1 \ident !B$, નહીં તો
ડાબો સિક્વન્ટ સ્વયંસિદ્ધ નથી;
અને $!C_2 \ident !A$ તથા $!E \ident !B$,
નહીં તો જમણો સિક્વન્ટ સ્વયંસિદ્ધ નથી.
આપણે સાબિતી સંપૂર્ણપણે પુનઃપ્રાપ્ત કરી છે:
  \[
  \Axiom$x : !A\land !B \fCenter x: !A\land !B$
  \RightLabel{\Elim\land[2]}
  \UnaryInf$x : !A\land !B \fCenter \proj{2}{x} : !B$
  \Axiom$x : !A\land !B \fCenter x: !A\land !B$
  \RightLabel{\Elim\land[1]}
  \UnaryInf$x : !A\land !B \fCenter \proj{1}{x} : !A$
  \RightLabel{\Intro\land}
  \BinaryInf$x : !A\land !B \fCenter \pair{\proj{2}{x}}{\proj{1}{x}} : 
    !B \land !A$
  \RightLabel{\Intro\lif}
  \UnaryInf$\fCenter \lambd[\typeof{x}{!A \land
  !B}][\pair{\proj{2}{x}}{\proj{1}{x}}] : (!A \land !B) \lif (!B \land !A)$
  \DisplayProof
  \]


\end{document}

